\documentclass[11pt]{article}

% Change "review" to "final" to generate the final (sometimes called camera-ready) version.
% Change to "preprint" to generate a non-anonymous version with page numbers.
\usepackage[review]{acl}

% Standard package includes
\usepackage{times}
\usepackage{latexsym}


% For proper rendering and hyphenation of words containing Latin characters (including in bib files)
\usepackage[T1]{fontenc}
% For Vietnamese characters
% \usepackage[T5]{fontenc}
% See https://www.latex-project.org/help/documentation/encguide.pdf for other character sets

% This assumes your files are encoded as UTF8
\usepackage[utf8]{inputenc}

% This is not strictly necessary, and may be commented out,
% but it will improve the layout of the manuscript,
% and will typically save some space.
\usepackage{microtype}

% This is also not strictly necessary, and may be commented out.
% However, it will improve the aesthetics of text in
% the typewriter font.
\usepackage{inconsolata}

%Including images in your LaTeX document requires adding
%additional package(s)
\usepackage{graphicx}

\usepackage{booktabs}
\usepackage{float}

% If the title and author information does not fit in the area allocated, uncomment the following
%
%\setlength\titlebox{<dim>}
%
% and set <dim> to something 5cm or larger.

\title{Decontaminated by Whose Standard? \\
A Label-Free Evaluation of Decontamination Filters}

% Author information can be set in various styles:
% For several authors from the same institution:
% \author{Author 1 \and ... \and Author n \\
%         Address line \\ ... \\ Address line}
% if the names do not fit well on one line use
%         Author 1 \\ {\bf Author 2} \\ ... \\ {\bf Author n} \\
% For authors from different institutions:
% \author{Author 1 \\ Address line \\  ... \\ Address line
%         \And  ... \And
%         Author n \\ Address line \\ ... \\ Address line}
% To start a separate ``row'' of authors use \AND, as in
% \author{Author 1 \\ Address line \\  ... \\ Address line
%         \AND
%         Author 2 \\ Address line \\ ... \\ Address line \And
%         Author 3 \\ Address line \\ ... \\ Address line}

\author{First Author \\
  Affiliation / Address line 1 \\
  Affiliation / Address line 2 \\
  Affiliation / Address line 3 \\
  \texttt{email@domain} \\\And
  Second Author \\
  Affiliation / Address line 1 \\
  Affiliation / Address line 2 \\
  Affiliation / Address line 3 \\
  \texttt{email@domain} \\}

%\author{
%  \textbf{First Author\textsuperscript{1}},
%  \textbf{Second Author\textsuperscript{1,2}},
%  \textbf{Third T. Author\textsuperscript{1}},
%  \textbf{Fourth Author\textsuperscript{1}},
%\\
%  \textbf{Fifth Author\textsuperscript{1,2}},
%  \textbf{Sixth Author\textsuperscript{1}},
%  \textbf{Seventh Author\textsuperscript{1}},
%  \textbf{Eighth Author \textsuperscript{1,2,3,4}},
%\\
%  \textbf{Ninth Author\textsuperscript{1}},
%  \textbf{Tenth Author\textsuperscript{1}},
%  \textbf{Eleventh E. Author\textsuperscript{1,2,3,4,5}},
%  \textbf{Twelfth Author\textsuperscript{1}},
%\\
%  \textbf{Thirteenth Author\textsuperscript{3}},
%  \textbf{Fourteenth F. Author\textsuperscript{2,4}},
%  \textbf{Fifteenth Author\textsuperscript{1}},
%  \textbf{Sixteenth Author\textsuperscript{1}},
%\\
%  \textbf{Seventeenth S. Author\textsuperscript{4,5}},
%  \textbf{Eighteenth Author\textsuperscript{3,4}},
%  \textbf{Nineteenth N. Author\textsuperscript{2,5}},
%  \textbf{Twentieth Author\textsuperscript{1}}
%\\
%\\
%  \textsuperscript{1}Affiliation 1,
%  \textsuperscript{2}Affiliation 2,
%  \textsuperscript{3}Affiliation 3,
%  \textsuperscript{4}Affiliation 4,
%  \textsuperscript{5}Affiliation 5
%\\
%  \small{
%    \textbf{Correspondence:} \href{mailto:email@domain}{email@domain}
%  }
%}

\begin{document}
\maketitle
\begin{abstract}
There is growing agreement that existing decontamination methods have limitations. Still, open post-training projects regularly report that their data has been decontaminated against public benchmarks. But do these filters even agree on what counts as contamination? We compare five decontamination filters from open datasets, applying them to the same 100,000 training documents and 502 benchmark questions. The results differ substantially, as filters flag anywhere from 0 to 141 questions.

Even though the filters disagree, this does not necessarily mean that any of them are wrong. To evaluate flagged items without manual annotation, we use benchmark questions published after the training corpus was collected. Since these questions did not yet exist, they could not have appeared in the corpus, so any flags on them must be false positives. The two filters that flag the most questions have false-positive rates of 12.9\% and 11.1\% on Dolci, and between 5\% and 16\% on two other corpora. We also find that common mathematical wording can cause false positives, which become more frequent as the corpus grows. Finally, we use the same controls to test how changing filter settings affects false-positive rates.
\end{abstract}

\section{Introduction}

Recent research has highlighted that benchmark contamination can inflate the performance of language models on evaluation tasks, undermining the reliability of benchmark-based comparisons \citep{chen-etal-2025-benchmarking-large,oren2024,sainz2023}. This has motivated efforts to develop more rigorous methods for detecting and preventing contamination in training datasets \citep{xu-etal-2025-dcr,controllable,detectable}. The underlying premise is that a benchmark score provides evidence of a model's capabilities only when the benchmark is held out from training. Papers introducing open-weight models commonly report efforts to mitigate contamination with respect to the benchmarks used for evaluation \citep{s1,tulu3,openthoughts,olmo2025olmo3}. However, the extent to which these claims are independently verified remains unclear.

Making matters worse, researchers and downstream leaderboard operators often treat decontamination claims as comparable \citep{leaderboardillusion}. However, two decontamination procedures are comparable only if they provide similar levels of assurance regarding what they detect, what they remove, and how they establish that a dataset is uncontaminated.

We evaluate five decontamination filters distributed with publicly released datasets -- s1 \citep{s1}, OpenThoughts-114K and OpenThoughts3 \citep{openthoughts}, Olmo~3's decon \citep{olmo2025olmo3}, and T\"{u}lu~3 \citep{tulu3} -- reimplementing each from its released source or running the released binary where available. The filters differ substantially in their detection mechanisms. We apply all five filters to a shared corpus of 100,000 training documents from Dolci \citep{olmo2025olmo3} against a benchmark suite of 502 items drawn from OlymMATH \citep{olymmath} and MathArena \citep{matharena}. The filters flag 127, 0, 141, 0, and 2 items, respectively. As a result, a dataset described as decontaminated under one procedure may not be considered decontaminated under another.

Showing that filters disagree does not establish that they are wrong. Assessing whether a flagged example is genuinely contaminated could rely on human annotation, but that approach is expensive and unreliable at exactly the ambiguous cases that decide it. We propose a better approach, by constructing a set of benchmark items for which the answer is already known. Competition mathematics publishes new problems continuously with known dates, and an item published after a training corpus was assembled cannot have been in it. Any filter that flags such an item therefore produces a false positive, without requiring human judgment about whether the match is considered contamination. This claim does not depend on the exact cutoff date, since no string filter can match an item that did not yet exist. While recent work has explored temporal signals in model performance to detect contamination \citep{ICLR2024_2d04d975}, these signals can be confounded by the structure of benchmark questions \citep{zhang-etal-2026-test}. Our approach instead examines the training data directly, using publication dates to identify items that could not have been present in the corpus.

We test the assumption behind this control rather than taking it for granted (Section~\ref{sec:control-assumption}). For the two filters that flag the most items, we measure false-positive rates of $12.9\%$ and $11.1\%$ on the primary corpus (Section~\ref{sec:false-positive-rates}), and between $5\%$ and $16\%$ on two further corpora, with the proportion remaining stable across corpus sizes (Section~\ref{sec:robustness}). We also identify a set of mechanistic errors caused by a small number of recurring phrases in competition mathematics (Section~\ref{sec:mechanisms}). Finally, we introduce a tuning procedure showing that n-gram size affects sensitivity on the control benchmarks rather than precision (Section~\ref{sec:calibration}).


\section{Related Work}

A significant amount of research on contamination has focused on examining training data. An alternative approach is to examine the model's behavior for evidence of contamination. Model-centered methods use membership-inference-style token-probability statistics \citep{shi2024detectingpretrainingdatalarge} or exchangeability arguments \citep{oren2024}. \citet{gemreview} conduct a systematic analysis of 55 papers and find that none of the methods is effective in all cases of contamination. Similarly, \citet{reliabilitygap} find that only 60\% of 335 tests using three leading methods produce correct results. \citet{detectable} explain this from an information-theoretic perspective, showing that detectability depends on both prevalence and sample size. Their results suggest that audit-scale samples may be one to two orders of magnitude too small to reliably detect contamination. Contamination detection in reasoning models is also fragile \citep{fragility}, while detectors that rely on LLM judges face contamination issues of their own \citep{preferenceleakage,judgerobustness}.

When the training corpus used by a model is available, overlap can be inspected directly. Another limitation is that n-gram-based filters detect only surface-level matches and can therefore miss paraphrased contamination \citep{yang2023rephrase}. Trying to answer how contaminated a corpus is might be the wrong approach. Instead, we ask whether contamination has a single answer in the first place, given that different procedures used to measure it disagree. Recent work on statistical guarantees for contamination detection \citep{controllable} assumes that there is a single true contamination rate that can be estimated reliably. Rather than trying to estimate this rate, we test whether that assumption holds in the first place.

In parallel, other work has focused on the reliability of audit pipelines. \citet{auditingtheaudit} use perturbations to examine the validity of these pipelines and identify five categories of failure modes. \citet{reliabilitygap} study model-side contamination detection under distribution shift and with small sample sizes, uncovering substantial limitations.  However, neither study directly compares published decontamination pipelines on a common corpus or reports how much their results differ.

\section{Filters and Experimental Setup}
\label{sec:filters}

The five filters differ in what they consider contamination. Some flag any shared n-gram, while others use whole-string similarity or multi-stage matching pipelines. To compare their behavior fairly, all filters were applied to the same training corpus and benchmark suite, using their released implementations rather than a shared implementation that could change how they work.

\subsection{Filter Implementations}
\label{sec:filter-implementations}
\begin{figure*}[t]
  \centering
  \includegraphics[width=\textwidth]{Abbildung_Corpus_Cutoff.jpg}
  \caption{Overview of the training corpus and benchmark sets, showing their composition and publication dates relative to the corpus cutoff.}
  \label{fig:corpus_overview}
\end{figure*}

\textbf{The filters use different procedures to identify potential contamination.} The n-gram matching introduced by \citet{gpt3} flags benchmark items based on shared n-grams in a single stage approach. Fuzzy string matching instead identifies matches whose whole-string similarity exceeds a predefined threshold. Decon uses a two-stage pipeline, where it first retrieves candidates based on 5-gram overlap, then expands matching clusters and retains only candidates with weighted scores above 0.8. Tülu 3 \citep{tulu3} uses Elasticsearch-based retrieval to identify candidate matches, making its detection coverage dependent on the retrieval stage.

\textbf{The largest difference between the single-stage filters lies in their tokenization.} s1 \citep{s1} splits text on whitespace after lowercasing, whereas OpenThoughts3 \citep{openthoughts} uses the Qwen2 tokenizer \citep{qwen2}. As a result, an 8-gram may span eight words in s1 but a different number of words in OpenThoughts3, since Qwen2 can split individual words into multiple tokens.

\textbf{Two filters require clarification because their implementations differ across sources.} The OpenThoughts repository \citep{openthoughts} contains two releases that use different decontamination methods. OpenThoughts-114K uses \texttt{fuzz.ratio} on raw strings with a cutoff of 95, while the later OpenThoughts3 marks an item contaminated by either a shared 13-gram or a \texttt{fuzz.ratio} above 75. We ran each filter against the release it was designed for. T\"{u}lu~3's \texttt{-{}-ngram\_size} and \texttt{-{}-match\_threshold} arguments have no default values in the released code. The paper specifies an 8-gram size and a match threshold of greater than 50\%, so we used those values and identify them as derived from the paper rather than the code \citep{tulu3}.

\subsection{Corpora and Benchmark Sets}
\label{sec:corpora}

Figure~\ref{fig:corpus_overview} shows the training corpus and the two benchmark sets used in our experiments, along with their composition and publication dates. The measurement set contains problems published before the corpus cutoff, meaning they could have been in the training data. The control set contains problems published well after the cutoff, so they could not have been included.

\textbf{Corpus (2,253,684 examples).} We use AllenAI's \texttt{Dolci-Think-SFT-32B} post-training dataset, developed for OLMo 3 \citep{olmo2025olmo3}. Since OLMo 3 was announced on November 20, 2025, we use November 2025 as an approximate cutoff date for the temporal negative control. We evaluate corpus sizes ranging from 12,500 to 390,000 examples to examine how flagging rates change with corpus size. Each Dolci example consists of a conversation between a user and an assistant, including the assistant's reasoning trace. We therefore check for benchmark problems in the user messages, where the original questions would appear if they were present in the data. This choice also follows the released implementations of the decontamination filters, which search the user-message field.

\textbf{Measurement set (502 items).} Our measurement set combines 200 items from OlymMATH \citep{olymmath} with problems from ten MathArena competitions held before November 2025 \citep{matharena}. These competitions are Apex 2025, the Apex shortlist, SMT 2025, CMIMC 2025, BRUMO 2025, AIME 2025, HMMT February 2025, both AIME 2024 contests, and HMMT February 2024. These problems were published before the corpus cutoff, so a filter may flag them because of genuine contamination or simply because of a matching phrase. The control set helps us distinguish between the two. 

\textbf{Control set (280 items).} The control set comprises eight MathArena competition sets published between January and June 2026, including AIME 2026, HMMT February 2026, USAMO 2026, three ArXivMath sets and two BrokenArXiv sets. To ensure a sufficient temporal separation from the training corpus, all items were published at least two months after the November 2025 cutoff.


\section{Measuring Filter Error with Temporal Controls}
\label{sec:temporal-controls}

Determining whether a flagged item is genuinely contaminated can require extensive manual inspection, making systematic verification impractical at scale \citep{golchin-surdeanu-2025-data}. Instead of manually checking each match, we use the corpus collection date to identify false positives. Since the corpus was assembled at time $t$, a benchmark item published after $t$ could not have been present in it. Any such item flagged by a decontamination filter therefore counts as a false positive.

Previous work has also explored temporal analysis as a way to detect contamination by comparing model performance on questions published before and after a model's knowledge cutoff \citep{ICLR2024_2d04d975, zhang-etal-2026-test}. However, these approaches examine changes in model performance, while we use publication dates to directly measure the false-positive rates of decontamination filters.

\subsection{Validating the Control Assumption}
\label{sec:control-assumption}

\begin{figure}[t]
  \centering
  \includegraphics[width=\linewidth]{Abbildung_False_Positive.jpg}
  \caption{Example of a false positive where the s1 filter flags a post-cutoff problem because it shares a phrase with a Dolci training example.}
  \label{fig:false_positive_example}
\end{figure}

Our false-positive analysis relies on the assumption that the control items were published after the training corpus was collected. This assumption requires particular care because Dolci combines data from multiple sources that may have been collected at different times \citep{olmo2025olmo3}. If any of these sources included data added after our assumed cutoff, the temporal control would be invalid.

We first inspect the publication dates of the control items. Our code confirms that every control item comes from a competition held after the training data was collected. We then manually inspect the five longest matches to determine whether they contain any of the original problem content. The longest match \textit{``m and n are relatively prime positive integers, find m + n''} is simply the ending of a phrase of an AIME problem. (Figure~\ref{fig:false_positive_example}) Three of the next four longest matches involve the same problem. The remaining match reads the same from \textit{``left to right as from right to left''}, which is only the description of a palindrome. These examples suggest that the filters can flag generic mathematical phrasing even when the shared text does not indicate that the benchmark problem itself appeared in the training data.

\begin{table}[t]
  \centering
  \small
  \setlength{\tabcolsep}{5pt}
  \begin{tabular}{@{}l rrrrr@{}}
    \toprule
    & \multicolumn{5}{c}{Longest shared span (words)} \\
    \cmidrule(l){2-6}
    & $<8$ & 8--9 & 10--12 & 13--14 & $\geq 15$ \\
    \midrule
    Measurement (502) & 375 & 91 & 29 & 7 & 0 \\
    Control (280)     & 244 & 25 &  9 & 1 & 1 \\
    \bottomrule
  \end{tabular}
  \caption{Longest word $n$-gram each benchmark item shares with any
  document in the corpus. Each item is counted once, in the bin of its
  longest match.}
  \label{tab:longest-match}
\end{table}

Table~\ref{tab:longest-match} summarizes the longest shared match for each benchmark item. Across the two sets, 127 measurement items and 36 control items share at least eight consecutive words with the corpus. This is consistent with the behavior of the s1 filter, which flags items based on 8-gram overlap \citep{s1}. The presence of these matches in the control set illustrates why shared text alone does not establish contamination.

\section{Disagreement and False Positives}
\label{sec:disagreement}

\begin{table}[t]
  \centering
  \small
  \setlength{\tabcolsep}{3pt}
  \begin{tabular}{@{}ll rr@{}}
    \toprule
    Filter & Rule & Measurement & Control \\
           &      & (502)       & (280)   \\
    \midrule
    decon             & two-stage &   0 (0.0)  &  0 (0.0)  \\
    OT-114K & fuzzy     &   0 (0.0)  &  0 (0.0)  \\
    T\"ulu~3          & coverage  &   2 (0.4\%)  &  0 (0.0)  \\
    s1                & 8-gram    & 127 (25.3\%) & 36 (12.9\%) \\
    OT3     & 13-gram / fuzzy & 141 (28.1\%) & 31 (11.1\%) \\
    \bottomrule
  \end{tabular}
  \caption{Benchmark items flagged by each filter on 100,000 Dolci
  documents, with percentages in parentheses. Control items cannot be
  in the corpus, so every control flag is a false positive.}
  \label{tab:filter-flags}
\end{table}

As observed in Table~\ref{tab:filter-flags}, the five filters applied to the same corpus and benchmarks reach very different conclusions. They separate into two groups with nothing in between. Three flag almost nothing, while two flag a quarter of the benchmark. There is no shared convention on where exactly contamination should sit. This section aims to answer that question. Are the filters finding different types of contamination? Or are the filters applying different standards for what counts as a match? The control set then tells us how many of those flags are wrong.

\subsection{Filter Disagreement}
\label{sec:filter-disagreement}

If the filters were finding the same contaminated items, we would expect their results to overlap. Instead, they often flag different items. Of the $203$ benchmark items flagged by at least one filter, $137$ ($67.5\%$; $95\%$ bootstrap CI $[61.1, 73.8]$) are flagged by exactly one filter. No item is flagged by more than three of the five filters. The highest pairwise Jaccard similarity is only $0.320$, meaning that even the two filters with the most overlap share less than a third of the items they flag. As we show in Section~\ref{sec:calibration}, changing the $n$-gram length also changes which items are flagged, but does not tell us whether those flags are correct. The amount of contamination a filter reports therefore depends heavily on which filter is used.

\texttt{decon}'s zero makes a sharper version of the same point. Our first
implementation reproduced its 5-gram candidate-generation stage but not the
cluster scoring applied on top; this version flagged $199$ of $200$ items where the released binary flags none. \footnote{We provide benchmark answers but not training-document answers, limiting question-only matches to $0.75$, below decon's $0.8$ threshold. Disabling answer indexing and rerunning the filter still produced no flags.} Candidate generation finds potential matches, not final ones, and treating the two as equivalent inverts the result. That zero describes what \texttt{decon} reports on this corpus, not how the corpus was filtered: Olmo~3 documents \texttt{decon} under midtraining, and its post-training data as filtered with the T\"{u}lu~3 toolkit at $8$-grams and a $0.5$ threshold \citep{olmo2025olmo3}.

\subsection{Measured False-Positive Rates}
\label{sec:false-positive-rates}

There are two core metrics we compare in our analysis. First, the \textbf{false-positive rate}, which is the number of flagged items in the control group divided by the total number of control items. Second, the \textbf{estimated false-discovery share}, which is the control rate divided by the measurement rate.

Both tested n-gram filters flag items that cannot have been in the corpus. The measured \textbf{false-positive rate} is 12.9\% for s1 (95\% CI [9.3, 16.8]) and 11.1\% for OpenThoughts3 (95\% CI [7.5, 15.0]), compared to measurement rates of 25.3\% (95\% CI [21.5, 29.1]) and 28.1\% (95\% CI [24.1, 32.1]), respectively. We calculate the confidence intervals by bootstrapping the benchmark items 10,000 times. Since the control set contains 280 items, the false-positive rate is also subject to uncertainty due to the limited number of control items.

The \textbf{estimated false-discovery share} is \(12.9/25.3 = 0.51\) for s1 and \(11.1/28.1 = 0.39\) for OpenThoughts3. This means that at most half of s1's output and at most 39\% of OpenThoughts3's is noise. Reading the ratio as an upper bound assumes that control questions and uncontaminated measurement questions attract spurious matches equally often. We test that assumption in Section~\ref{sec:matched} rather than taking it for granted. Subtracting the control rate from the measurement rate leaves 12.4 percentage points for s1 and 17.0 percentage points for OpenThoughts3.

\section{Robustness Across Corpora and Scales}
\label{sec:robustness}

Since false positives are common in the control set, we investigate whether their frequency increases as the corpus grows.

\subsection{Scaling with Corpus Size}
\label{sec:scaling}
 
\begin{figure*}[t]
  \centering
  \includegraphics[width=\textwidth]{scaling_figure.png}
  \caption{Flag rates of s1 and OpenThoughts3 across Dolci corpus sizes from 12.5K to 390K documents. (a) Both filters flag more measurement and control items as the corpus grows. (b) The ratio of control to measurement flag rate, an upper bound on the share of each filter's output that is noise.}
  \label{fig:scaling}
\end{figure*}

Previous work showed that measured contamination depends on the matching criterion used \citep{singh2024evaluationdatacontaminationllms}, but does not address whether false positives become more frequent as the training corpus grows.
To investigate this, we gradually increased the Dolci corpus from 12,500 to 390,000 records, to test the 502 measurement items and 280 control at seven corpus sizes. Because of higher flag rates, we focus on s1's 8-gram and OpenThoughts3's 13-gram. Figure~\ref{fig:scaling} shows the scaling trends, and Table~\ref{tab:scaling} in Appendix~\ref{app:scaling} reports the complete rates.

Both filters flag more measurement and control items as the corpus grows. For s1, the percentage of flagged measurement items increases from 13.7\% to 32.9\%, while the control false-positive rate rises from 6.8\% to 17.9\%. We observe a similar pattern for OpenThoughts3's Qwen tokenizer-based 13-gram filter, where the measurement flag rate increases from 12.9\% to 37.6\% and the control false-positive rate from 5.4\% to 17.1\%.

Searching more training records increases the chance of finding matching text, which is why corpus size matters when comparing reported flag rates.

\subsection{Replication Across Corpora}
\label{sec:replication}

\begin{table}[t]
  \centering
  \small
  \begin{tabular}{@{}llrr@{}}
    \toprule
    Corpus & Filter & Measure & Control \\
    \midrule
    T\"ulu~3    & s1 8-gram    & 3.3\%  & 6.2\%  \\
    (Nov 2024)  & OT3 13-gram  & 1.7\%  & 6.9\%  \\
                & fuzzy / cov  & 0.0\%  & 0.0\%  \\
    \midrule
    Nemotron    & s1 8-gram    & 15.4\% & 5.4\%  \\
    (May 2026)  & OT3 13-gram  & 24.1\% & 16.1\% \\
                & fuzzy / cov  & 0.1\%  & 0.0\%  \\
    \bottomrule
  \end{tabular}
  \caption{Two further corpora, $100{,}000$ documents each. T\"ulu~3 has $60$ measurement and $722$ control items, Nemotron $680$ and $261$. Each split is derived from that corpus's own release date.}
  \label{tab:corpora}
\end{table}

We repeat the comparison for four filters on 100,000 records each from the T\"ulu~3 mathematics subset and Nemotron-SFT-Math-v4 \citep{nemotron_sft_math_v4}. Each corpus uses a temporal split based on its own release date. The T\"ulu corpus has 60 measurement entries and 722 controls, while Nemotron has 680 measurement entries and 261 controls. Decon was excluded from these additional experiments.

As shown in Table~\ref{tab:corpora}, s1 incorrectly flags 45 controls (6.2\%) on T\"ulu~3, while OpenThoughts3 flags 50 (6.9\%). On Nemotron, s1 flags 14 controls (5.4\%) and OpenThoughts3 flags 42 (16.1\%). OpenThoughts3 also flags 164 measurement items (24.1\%) on Nemotron. Whole-string fuzzy matching and the coverage filter produce no control flags in either corpus. These results suggest that false positives from n-gram matching are not specific to Dolci, even though their frequency varies across corpora and filters.

Which filter produces more false positives also depends on the corpus. On Dolci, s1 has a higher false-positive rate, while on T\"ulu~3 and Nemotron, OpenThoughts3 has a higher rate. However, the benchmark questions also differ between corpora. Because Nemotron has a later cutoff date, its control set only contains arXiv-derived questions. This makes it difficult to tell whether differences in false-positive rates come from the training corpus or the control questions. But we can still compare the filters within each corpus, since they are all tested on the same questions.

\subsection{Comparing the Same Competitions Across Years}
\label{sec:matched}

The bound we set in \S\ref{sec:false-positive-rates} assumes that measurement and control questions are equally likely to produce false positives. Comparing 2025 and 2026 editions of AIME and HMMT helped us check this. Both competitions follow similar problem-writing conventions each year, but only the 2025 questions could have appeared in the training corpus. The 2026 questions are too new, so any matches with them must be false positives. If the filters flag both years at similar rates, this suggests that many of the 2025 matches are also false positives rather than actual contamination.

\begin{table}[t]
  \centering
  \small
  \begin{tabular}{@{}llrr@{}}
    \toprule
    Competition & Group & s1 & OT3 \\
    \midrule
    AIME 2025      & measure & 60.0\% & 53.3\% \\
    AIME 2026      & control & 56.7\% & 36.7\% \\
    HMMT Feb 2025  & measure & 20.0\% & 26.7\% \\
    HMMT Feb 2026  & control & 21.2\% & 12.1\% \\
    \bottomrule
  \end{tabular}
  \caption{Flag rates for the same competitions one year apart, using $100{,}000$ Dolci records. Only the 2025 editions could appear in the corpus. OT3 refers to the Qwen 13-gram filter.}
  \label{tab:matched}
\end{table}

s1 shows almost the same rate across both years. On AIME, it flags 60.0\% of the 2025 questions and 56.7\% of the 2026 questions. On HMMT, it even flags slightly more questions from 2026 than from 2025 (21.2\% versus 20.0\%). Since the 2026 questions could not have been in the corpus, this suggests that much of s1's flagging comes from common wording rather than contamination. The similarity of these rates supports the assumption behind the bound in Section~\ref{sec:false-positive-rates}, at least for these two competitions. OpenThoughts3 shows a larger gap between the years on both competitions, which could mean that some of its 2025 matches reflect actual contamination.

Nevertheless, we only consider two of the ten measurement competitions, and AIME questions tend to follow similar patterns. We therefore cannot conclude that this holds for all measurement benchmarks, although our results suggest that it may.

\section{Error Analysis and Control-Based Calibration}
\label{sec:error-analysis}

The flag counts do not explain why the filters produce false positives in the first place. We therefore examine the matched text more closely and test how changes to the filter settings affect the results.

\subsection{Where False Positives Come From}
\label{sec:mechanisms}

Sliding-window-based filters can flag text simply because it shares certain strings with a benchmark question. The problem is that matching a few strings does not necessarily mean the two problems are asking the same thing.

An example from s1 matches two questions beginning with \textit{``What is the smallest positive integer $n$ such that''}. One question concerns common complex roots of two polynomials, while the other concerns a terminating decimal. Many of the flags are caused by formatting rather than meaningful overlap between problems. For example, a control question about an isosceles triangle is matched with a training question about a hyperbola simply because both contain \texttt{points \textbackslash( P \textbackslash) and \textbackslash( Q \textbackslash)}. The formatting makes this phrase count as eight separate tokens, even though the two questions have little in common.

The way mathematical expressions are formatted can also affect text matching. Previous work has addressed inconsistencies in mathematical notation by standardizing expressions into a common LaTeX format \citep{ICLR2026_84f421a2}. But formatting can still lead to false positives in n-gram matching.

To better understand why the filters flag certain matches, we used an LLM judge to label 565 distinct candidate pairs, without revealing filter identities or metadata about the temporal controls. We collected up to three saved candidates per retrieval branch and benchmark entry. Of the 315 pairs retrieved by s1, 171 (54.3\%) were labelled generic overlap and 34 (10.8\%) unrelated. For the 325 pairs retrieved by the Qwen 13-gram component, these numbers were 89 (27.4\%) and 108 (33.2\%), respectively. This suggests that s1 more often matches common problem wording, while the Qwen component more frequently matches unrelated problems. Looking specifically at the controls, the judge labelled 50 of 114 pairs as generic overlap, 28 as related but distinct, and 36 as unrelated. While the temporal split tells us which matches are false positives, the LLM annotations help explain why those matches occur. Appendix~\ref{app:mechanism} traces all 36 of s1's control false positives to 73 distinct 8-grams. Table~\ref{tab:judge-labels} gives the full label distribution. Appendix~\ref{app:judge} describes the judging procedure and provides additional examples.

\subsection{Tuning Without Labels}
\label{sec:calibration}

\begin{figure}[t]
  \centering
  \includegraphics[width=\linewidth]{Tuning wo Labels.png}
  \caption{Measurement against control flag rates as each filter's setting
  varies, traced without labels on 100K Dolci documents. Rings mark
  shipped defaults. The OpenThoughts3 ring covers its 13-gram component
  only, not the 28.1\% reported in Table~\ref{tab:filter-flags}. Points
  beyond 20\% false positives are omitted, since at $n{=}5$ both $n$-gram
  families flag almost every item.}
  \label{fig:frontier}
\end{figure}

With our temporal design of the experiments, we can examine how parameter choices affect false positives without manually labelling every match. We test how changing the n-gram length for the word-based and Qwen filters, as well as the similarity threshold for whole-string fuzzy matching, affects the results. Figure~\ref{fig:frontier} plots each filter's measurement flag rate against its control false-positive rate as its settings are varied. The rings mark each project's default, and they differ considerably. The complete parameter sweeps appear in Table~\ref{tab:threshold-sweep} in Appendix~\ref{app:tuning}.

Increasing the word $n$-gram length from eight to thirteen reduces the measurement flag rate from 25.3\% to 1.4\% and the control false-positive rate from 12.9\% to 0.7\%. This shows that requiring longer matches reduces false positives, but also means that far fewer older benchmark questions are flagged. The ratio between the two barely moves, from $0.51$ to $0.50$, so $n$-gram length changes how much a filter flags without changing what fraction of its output is wrong.

The same pattern happens with whole-string fuzzy matching when changing its similarity threshold. The default threshold of 95 does not flag any items. Lowering it to 70 increases the measurement flag rate to 2.8\%, with no false positives on the controls. At 65, the measurement flag rate rises to 6.2\%, while the control false-positive rate reaches 0.7\%. To compare, word 10-gram matching flags 7.2\% of measurement items and 3.9\% of controls.

\subsection{Phrase-Agreement Filtering}
\label{sec:phrase-agreement}

The examples in Section~\ref{sec:mechanisms} motivate testing whether agreement across several positions reduces matches driven by a single conventional phrase. We evaluate a question-only adaptation that samples up to twenty position-distributed 8-word windows from each benchmark question. It flags the question when at least a specified fraction of those windows matches within a single training document.

Table~\ref{tab:phrase-agreement} reports the results. Requiring 20\% agreement flags three measurement entries (0.6\%) and no controls. Reducing the requirement to 10\% increases these counts to 22 (4.4\%) and seven (2.5\%). At 5\%, the filter flags 87 measurement entries (17.3\%) and 22 controls (7.9\%).

\begin{table}[t]
\centering
\small
\setlength{\tabcolsep}{5pt}
\begin{tabular}{@{}lrr@{}}
\toprule
Required fraction & Measurement & Control \\
\midrule
$\geq 0.50$ & 0 (0.0)   & 0 (0.0) \\
$\geq 0.30$ & 2 (0.4)   & 0 (0.0) \\
$\geq 0.20$ & 3 (0.6)   & 0 (0.0) \\
$\geq 0.10$ & 22 (4.4)  & 7 (2.5) \\
$\geq 0.05$ & 87 (17.3) & 22 (7.9) \\
\bottomrule
\end{tabular}
\caption{Question-only phrase-agreement filtering on
100,000 Dolci records. Counts and percentages refer
to 502 measurement entries and 280 controls.
The filter samples up to 20 position-distributed
8-word windows and requires the specified fraction
to match within a single training document.}
\label{tab:phrase-agreement}
\end{table}

Stricter agreement requirements reduce the observed control false positives, alongside a reduction in measurement flags. Agreement across sampled positions does not eliminate errors at all settings: sampled windows may overlap or fall within the same conventional wording. Moreover, the low measurement yield at stricter settings leaves their ability to detect genuine contamination unresolved.

\section{Practical Implications for Decontamination}
\label{sec:practice}

The authors of the filters we tested do not appear to have chosen their default settings based on measured false-positive rates. The OpenThoughts-114K filter applies a fuzzy matching threshold of 95 and flags nothing at all on our corpus. Tuned down to 65, the same method flags 6.2\% of the measurement items with just 0.7\% false positives. A word 10-gram flags a similar share, 7.2\%, but with a much larger false-positive rate of 3.9\%.

By collecting benchmark questions published after their training corpus was assembled, new projects could evaluate their own filters. This would allow them to measure false-positive rates and use these results to choose a threshold instead of relying on the default. The chosen threshold should then be tested on a separate control set \citep{cawley2010over}. Projects should also report implementation details, such as the tokenizer and threshold used, since without these details it is difficult to know what a claim of decontamination actually means \citep{leaderboardillusion}.

\section*{Limitations}

A control flag is wrong because the question did not exist when the corpus was built. That part needs no further checking. The other direction is invisible to us. We have no items known to be contaminated, so we cannot estimate the amount of contamination or the effectiveness of each filter. We can also not conclude that fewer control flags indicate a better filter and how the decisions related to filtering affect model performance. Another limitation is that our control set contains 69 competition questions and 211 questions derived from research. In contrast, the measurement set consists almost entirely of competition questions. The two groups show very different false-positive rates. For s1, 40.6\% of competition controls are flagged compared to only 3.8\% of research-derived controls. For OpenThoughts3, the rates are 26.1\% and 6.2\%. As a result, our overall false-positive rates depend a lot on the composition of the control set and may underestimate the false-positive rates for competition questions (Appendix~\ref{app:strata}). 

The reported results are based on benchmark entries rather than distinct problems. The 502 measurement items contain 494 distinct question texts, because the Apex shortlist reuses problems that also appear in AIME, HMMT and BRUMO. We keep the items as they originally appeared in the benchmark, since projects using these filters would encounter the same duplicates. This has the consequence, that every duplicated item that is flagged contributes twice to the rate, which moves s1 from 25.3\% to 24.7\% and OpenThoughts3 from 28.1\% to 27.9\% on the deduplicated set.

The way we implemented Tülu's filters replaces the original Elasticsearch retrieval implementatin with exhaustive token matching. So its results describe an adapted mechanism to the released pipeline, and the LLM annotations have not been validated against human judgments.

The clearest next step is a controlled training experiment, measuring whether different filtering decisions change scores on benchmarks or remove useful training data. That would connect filter error to a score people care about which our experiments do not do.

\bibliography{custom}

\clearpage
\appendix

\section{Results Across Corpus Sizes}
\label{app:scaling}

Table~\ref{tab:scaling} reports the complete corpus-size sweep. All settings use the same 502 measurement entries and 280 controls. Each larger corpus prefix contains the smaller ones, so the rows are not independent samples. The Qwen results cover the 13-gram component alone, excluding fuzzy matching.

\begin{table}[!t]
\centering
\small
\setlength{\tabcolsep}{3pt}
\begin{tabular}{@{}rrrrr@{}}
\toprule
& \multicolumn{2}{c}{s1: word 8-gram} & \multicolumn{2}{c}{Qwen 13-gram} \\
\cmidrule(lr){2-3}\cmidrule(l){4-5}
Records & Meas. & Control & Meas. & Control \\
\midrule
12,500  & 13.7 & 6.8  & 12.9 & 5.4 \\
25,000  & 15.7 & 8.2  & 16.9 & 6.1 \\
50,000  & 18.1 & 11.1 & 21.9 & 8.2 \\
100,000 & 25.3 & 12.9 & 27.5 & 11.1 \\
150,000 & 27.1 & 14.3 & 32.9 & 13.2 \\
250,000 & 30.3 & 16.1 & 35.7 & 15.4 \\
390,000 & 32.9 & 17.9 & 37.6 & 17.1 \\
\bottomrule
\end{tabular}
\caption{Benchmark flag rates (\%) across nested Dolci prefixes. Measurement and control denominators are 502 and 280, respectively. Control flags are false positives under the temporal assumption.}
\label{tab:scaling}
\end{table}

\section{Candidate-Pair Analysis}
\label{app:judge}

We formed the judge pool from saved s1 and Qwen 13-gram matches, fuzzy candidates scoring at least 75, and candidates exceeding 50\% per-document coverage under the audit's whitespace tokenizer. Up to three candidates were saved per matching method and benchmark entry. Removing duplicate text pairs left 565 distinct pairs. The coverage candidates come from the audit adaptation, rather than the final spaCy implementation.

We used \texttt{gpt-6-astra} with high reasoning effort. The judge received only the two question texts, without filter identities, matching scores, or temporal-group metadata. Table~\ref{tab:judge-labels} gives the complete label distribution. The pool contains 451 measurement pairs and 114 control pairs. Pairs retrieved by more than one method were judged once.

\begin{table}[!t]
\centering
\small
\setlength{\tabcolsep}{4pt}
\begin{tabular}{@{}lrrr@{}}
\toprule
Label & Meas. & Control & Total \\
\midrule
Same problem         & 0   & 0   & 0 \\
Related but distinct & 187 & 28  & 215 \\
Generic overlap      & 164 & 50  & 214 \\
Unrelated            & 100 & 36  & 136 \\
Uncertain            & 0   & 0   & 0 \\
\midrule
Candidate pairs      & 451 & 114 & 565 \\
Distinct questions   & 198 & 51  & 249 \\
\bottomrule
\end{tabular}
\caption{LLM labels for the selected candidate pairs. The final row counts distinct benchmark question texts. These labels describe the candidate pool, not the presence or absence of contamination elsewhere in the corpus.}
\label{tab:judge-labels}
\end{table}

Among s1's 315 pairs, 171 were labelled generic overlap, 110 related but distinct, and 34 unrelated. Among the Qwen component's 325 pairs, the corresponding counts were 89, 128, and 108. The per-document coverage and fuzzy branches contributed 6 and 4 pairs respectively. These groups overlap because a pair can be retrieved by more than one method.

We repeated the judgments for 20 randomly selected pairs (seed 42), using the same model and prompt without showing the original labels. The relationship label agreed in 19 cases. This measures repeatability, not agreement with human judgments.

\paragraph{A shared definition.}
One AIME 2026 control asks about a floor operation applied to a series, while its matched training question asks which integers can be expressed as a sum of floor functions. The s1 match is the phrase ``the greatest integer less than or equal to''. The definition is shared, but the problems are different.

\paragraph{Shared numerical notation.}
An OlymMATH question about the distances of cube vertices from a plane matches a training question about drawing numbered balls. The Qwen 13-gram matches occur within number sequences and their \LaTeX{} delimiters. The judge labelled the pair unrelated.

\section{Complete Parameter Sweeps}
\label{app:tuning}

Table~\ref{tab:threshold-sweep} reports the complete sweeps on 100,000 Dolci records. Each setting uses 502 measurement entries and 280 controls. The blocks vary separate rules: the Qwen block excludes fuzzy matching, and its token counts are not equivalent to whitespace-token counts.

\begin{table*}[!t]
\centering
\small
\setlength{\tabcolsep}{6pt}
\begin{tabular}{@{}rrr rrr rrr@{}}
\toprule
\multicolumn{3}{c}{Word $n$-grams} & \multicolumn{3}{c}{Qwen $n$-grams} & \multicolumn{3}{c}{Whole-string fuzzy} \\
\cmidrule(lr){1-3}\cmidrule(lr){4-6}\cmidrule(l){7-9}
$n$ & Meas. & Control & $n$ & Meas. & Control & Threshold & Meas. & Control \\
\midrule
5  & 93.0 & 72.5 & 5  & 99.8 & 99.6 & 95 & 0.0  & 0.0 \\
8  & 25.3 & 12.9 & 8  & 86.3 & 68.2 & 90 & 0.0  & 0.0 \\
10 & 7.2  & 3.9  & 10 & 62.0 & 36.4 & 85 & 0.0  & 0.0 \\
13 & 1.4  & 0.7  & 13 & 27.5 & 11.1 & 80 & 0.0  & 0.0 \\
15 & 0.0  & 0.4  & 15 & 14.3 & 6.8  & 75 & 0.8  & 0.0 \\
20 & 0.0  & 0.0  & 20 & 4.2  & 2.5  & 70 & 2.8  & 0.0 \\
25 & 0.0  & 0.0  & 25 & 1.2  & 0.4  & 65 & 6.2  & 0.7 \\
30 & 0.0  & 0.0  & 30 & 0.4  & 0.0  & 60 & 15.3 & 4.6 \\
\bottomrule
\end{tabular}
\caption{Complete parameter sweeps on 100,000 Dolci records. Entries are flag rates (\%) over 502 measurement entries and 280 controls. Each block varies one matching rule independently.}
\label{tab:threshold-sweep}
\end{table*}

Longer matches and higher fuzzy thresholds generally reduce the number of flags. Measurement entries have no contamination labels, so their flag rates should not be interpreted as recall.

\section{Which Phrases Cause the False Positives}
\label{app:mechanism}

Every one of s1's 36 control false positives traces to one of just 73 distinct 8-grams, and none is a fragment of a leaked problem. They are the conventional phrasing of competition mathematics. Table~\ref{tab:mechanism} lists the most frequent, with the number of training documents each appears in.

\begin{table}[!t]
\centering
\small
\begin{tabular}{@{}lr@{}}
\toprule
Shared 8-gram & Docs \\
\midrule
find the sum of all possible values of     & 166 \\
$a$ and $b$ be positive integers such that &  67 \\
is not divisible by the square of any      &  67 \\
let $a$ and $b$ be positive integers such  &  67 \\
not divisible by the square of any prime.  &  42 \\
\bottomrule
\end{tabular}
\caption{The five most frequent 8-grams behind s1's control false positives, with the number of training documents containing each. All 36 false positives trace to 73 such phrases.}
\label{tab:mechanism}
\end{table}

AIME problems conventionally close by asking for $m + n$ where $m$ and $n$ are relatively prime, and ``not divisible by the square of any prime'' is standard framing for a radical-form answer. Every problem of the type shares these conventions, including problems written months after the corpus was assembled. An eight-word window can be filled entirely by convention while carrying no problem-specific content, which is why raising $n$ does not escape the problem. The phrases are themselves long.

\section{Control Set by Question Type}
\label{app:strata}

Table~\ref{tab:strata} gives the full breakdown of control false positives by question type. Only 69 of the 280 control items are competition problems; the remaining 211 are derived from research papers.

\begin{table}[!t]
\centering
\small
\begin{tabular}{@{}lrr@{}}
\toprule
Filter & Competition (69) & Research (211) \\
\midrule
s1 8-gram         & 28 (40.6\%) & 8 (3.8\%) \\
OpenThoughts3     & 18 (26.1\%) & 13 (6.2\%) \\
OpenThoughts-114K & 0 (0.0\%)   & 0 (0.0\%) \\
T\"ulu~3          & 0 (0.0\%)   & 0 (0.0\%) \\
\bottomrule
\end{tabular}
\caption{Control false positives split by question type. The pooled rates in Table~\ref{tab:filter-flags} are weighted averages over these two strata.}
\label{tab:strata}
\end{table}

The two filters that flag nothing on either stratum are unaffected by the composition of the control set. For the two that do fire, the pooled rate sits between the two stratum rates and moves with their relative sizes.

\end{document}
